\documentclass[12pt]{article}
\title{Digest of The New York Times}
\author{The New York Times}

% % \definecolor{dark-green}{rgb}{0,0.25,0}
% \usepackage[utf8]{inputenc}

\usepackage{fontspec,xunicode,xltxtra}
\setmainfont[Ligatures=Required,Mapping=tex-text]{Georgia}
% \setmainfont[Ligatures=Required,Mapping=tex-text]{Palatino Linotype}
\setsansfont[Ligatures=Required,Mapping=tex-text]{Myriad Pro}
\setmonofont[Ligatures=Required,Mapping=tex-text]{Courier New}
\usepackage{xeCJK}
\setCJKmainfont{宋体}
\setCJKmonofont{黑体}

% \usepackage{indentfirst}


\usepackage[dvips,dvipsnames,svgnames]{xcolor}
\definecolor{light-gray}{gray}{0.95}

%生成PDF的链接
\usepackage[dvipdfmx]{hyperref}
\hypersetup{
    bookmarks=true,         % show bookmarks bar?
    bookmarksopen=true,
    pdfpagemode=UseNone,    % options: UseNode, UseThumbs, UseOutlines, FullScreen
    pdfstartview=FitH,
    pdfborder=1,
    pdfhighlight=/P,
    pdfauthor={wuxch},
    unicode=true,           % non-Latin characters in Acrobat’s bookmarks
    colorlinks,             % false: boxed links; true: colored links
    linkcolor=blue,         % color of internal links
    citecolor=blue,        % color of links to bibliography
    filecolor=magenta,      % color of file links
    urlcolor=cyan           % color of external links
}
% 使用目录
\usepackage{makeidx}
% 使用图形
\usepackage{graphicx}


\usepackage{titletoc}
\dottedcontents{section}[1.5em]{}{2.3em}{1pc}

% % 页眉页脚的设定
% \usepackage{fancyhdr}
% \pagestyle{fancy}
% \lhead{\includegraphics[height=0.2in]{zte_logo}}

% \usepackage{footmisc}
\usepackage{pfnote}
\usepackage{fnpos}

% % 嵌入的代码显示
% \usepackage{listings}
% \lstset{language=C++, breaklines, extendedchars=false}
% \lstset{keywordstyle=\color{blue},
%         commentstyle=\color{SeaGreen},
%         stringstyle=\ttfamily,
%         showstringspaces=false,
%         backgroundcolor=\color{light-gray}}
% \lstset{basicstyle=\small}

\usepackage{booktabs}

\usepackage{tipa}
\usepackage{phonetic}
\usepackage[gen]{eurosym}

% \usepackage{CJK}
\usepackage{makeidx}
% \usepackage{glossary}

\usepackage{fancyhdr}

\usepackage{lastpage}
\cfoot{\textsf{\thepage\ of \pageref*{LastPage}}}

\pagestyle{fancy}
\fancyhead{}
\renewcommand{\headrulewidth}{0pt}
\pagenumbering{arabic}

\usepackage{graphicx}
\usepackage{picinpar}

\setlength{\parindent}{0in}
\setlength{\parskip}{2.5ex plus4mm minus3mm}
% \setlength{\parskip}{.3\baselineskip}

\setcounter{section}{0}

% bibtex
\usepackage[super]{natbib}
% \usepackage[globalcitecopy,labelstoglobalaux]{bibunits}
% \makeatletter
% \renewcommand{\@bibunitname}{\jobname.\the\@bibunitauxcnt}
% \makeatother

% \renewcommand{\refname}{缩略语}
% \renewcommand{\refname}{}

\makeatletter                   % Reference list option change
\renewcommand\@biblabel[1]{#1:  } % from [1] to 1: 
\makeatother                    %

\newcommand{\citenumfont}[1]{\textsf{#1}}

% \usepackage{tipa}
% \usepackage{phonetic}

\newcommand{\phonetic}[1]{{\fontspec[Colour=CC00CC]{Microsoft Sans Serif} #1}}


%%% Local Variables:
%%% TeX-master: "medical_trivia"
%%% End:

\usepackage{config}

\makeindex
\begin{document}
\date{}
\thispagestyle{empty}

\begin{figure}
\includegraphics*[width=0.7\textwidth]{The_New_York_Times_logo.png}
\vspace{-15ex}
\end{figure}
% \renewcommand\contentsname{\textsf{Digest of The New York Times}}
\renewcommand\contentsname{}
{\footnotesize\textsf{\tableofcontents}}

\clearpage
\setcounter{page}{1}

\section{Market Defies Fear of Real Estate Bubble in China}

\lettrine{T}{he}\mycalendar{Mar.'10}{05} spacious duplex comes with crocodile-skin bedposts,
hand-carved bronze doors inlaid with Swarovski crystals -- and a \$45 million price tag.

\myphoto{Market.Defies.Fear.of.Real.Estate.Bubble.in.China.jpg}{The interior of a duplex show
apartment at the Tomson Riviera, a luxury apartment complex located in Pudong, Shanghai.}

It is still on the market, but Charles Tong, the developer of Tomson Riviera, a luxury riverfront
complex in the heart of the financial district here, says he is having no trouble finding takers for
similarly priced units.

``We're selling three to four apartments every month,'' said Mr.~Tong, seated in a white Versace
easy chair. ``Now, people here want something more luxurious; they'd like a new lifestyle.''

Everyone agrees China is in the middle of a spectacular real estate boom. The question is whether it
is in the middle of a rapidly growing real estate bubble.

When other recent booms collapsed -- in the United States, for instance -- they depressed entire
economies. In China's case, a bursting bubble could affect much of the world. China is the
fastest-growing large economy and, so far, a main engine pulling the world out of recession.

Beijing is clearly concerned. Authorities have recently moved to rein in the easy credit that has
helped finance China's hyperdevelopment, including making it more difficult for home buyers to take
out a second mortgage.

Last year, a record \$560 billion of residential property was sold in China, an increase of 80
percent from the year before, according to government statistics that are widely considered
reliable. And with prices soaring, developers are scrambling to build more mansions, villas and
high-rise apartments with names like Rich Gate, Park Avenue and Palais de Fortune.

Signs of exuberance are everywhere. An investor in Shanghai recently bought 54 apartments in a
single day; a villa sold for \$30 million last year; and in December a consortium of developers paid
more than \$3.5 billion for a huge tract of land in Guangzhou, one of the highest prices paid for
any property, anywhere. In the city of Tianjin, in north China, developers have created a \$3
billion ``floating city,'' a series of islands built on a natural reservoir, featuring villas,
shopping malls, a water amusement park and what they say will be the world's largest indoor ski
resort.

``This is wild,'' said Andy Xie, a former Morgan Stanley economist who is now an independent
analyst. ``By all the traditional measures, like rental yield, this is a bubble.''

Speculators are snapping up properties on the expectation that prices will continue to rise, as
prices have nearly every year for more than a decade. And powerful developers are working with local
governments to transform old cities into urban dreamscapes.

But Shanghai, China's wealthiest and most dazzling city, is the epicenter of the boom. Prices here
have risen more than 150 percent since 2003, pushing the price of a typical 1,100-square-foot
apartment up to \$200,000, according to real estate experts. (Shanghai residents typically earn less
than \$5,000 a year.)

A buying frenzy has gripped the city, leading to billion-dollar land auctions and long waiting
lists.

``The speed you buy a house here is faster than you buy vegetables,'' said Andy Xiang, an
advertising executive who recently put down a large cash down payment to get the right to pay \$1.3
million for an apartment in the city's exclusive Xintiandi area.

Few residences, though, are as upscale as Tomson Riviera, which consists of four golden-hued towers
overlooking the Huangpu River, with a central garden mapped out in the shape of a dragon. The
apartment complex's entrance has original artworks by Salvador Dalí and well-known Chinese artists.
The apartments, a few of which have been decorated by Armani and Fendi, as well as Versace, lease
for \$7,000 to \$17,000 a month -- to high-level executives from companies like General Motors.

Those who buy an apartment here tend to be extremely wealthy, like Liu Yiqian, an eccentric Shanghai
entrepreneur whom Forbes magazine says is worth about \$540 million.

Mr.~Liu, 47, got his start driving a taxicab in Shanghai but eventually made a fortune investing in
the stock market. In an interview this week, he acknowledged owning hundreds of apartments in
Shanghai (he said he could not remember exactly how many), including a 6,000-square-foot apartment
in Tomson Riviera, which he bought in 2008 for about \$11.5 million.

``I invest in properties,'' Mr.~Liu said, noting that he also collects art, antiques and jade. ``I
think in Shanghai in five to seven years the real estate prices will be even higher.''

As they try to modulate the market, local and central governments here are walking a thin line. Land
sales were a major source of government revenue, raising about \$234 billion last year, an amount
equal to over a third of the cost of China's half-trillion-dollar stimulus program.

Whether the country is in the middle of a bubble has become the subject of a debate. Some
economists, like Nicholas R.~Lardy at the Peterson Institute for International Economics in
Washington, say the housing boom is being propelled by a huge urbanization push that is creating
premium-priced houses.

Other analysts say prices are being propped up by greedy developers and government policies that are
making housing increasingly unaffordable for the masses migrating to big cities.

Despite the fear of a bubble here, Mr.~Tong said his prices were just right, particularly because of
so much hidden wealth in China. The publicly listed company is controlled by his family.

``I have a friend,'' he said. ``She makes maternity clothes. Her company has 20 percent of the
world's market share, and they're not even a listed company.''

Still, Tomson's prices are soaring. The most recent apartment sold for about \$2,300 a square foot.
The average luxury apartment in Manhattan sold for just under \$1,900 a square foot in the fourth
quarter of 2009, according to Prudential Douglas Elliman real estate.

Indeed, for the price of a Tomson apartment in Shanghai, a buyer could easily purchase a
6,000-square-foot home in Los Angeles built by Frank Lloyd Wright and now for sale (\$10.5 million),
or a 52-acre site with a 22-room residence in New Canaan, Conn. (\$24 million).

But a sales agent at Tomson Riviera says this is the future financial capital of the world, not the
dying one.

``Look at this bronze door,'' said Wang Yaodong. ``That costs \$50,000! Look at these Gaggenau
appliances. They were made in Germany.'' The glasses were imported from Belgium, the Jacuzzi from
Italy. And don't worry about losing your key, he said, ``This lock can read the palm of your hand.''

\section{For Sochi Olympics and Beyond: Choose Sites Wisely}

\lettrine{W}{ait}\mycalendar{Mar.'10}{05} out the post-Olympic rush for the exits, and it is
striking how quickly the normal rhythms and routines take over.

In Vancouver this week, less than 48 hours after the closing ceremony, the only thing I came across
during a long walk that was wearing bright red and a huge maple leaf was a seaplane coming in for a
landing in the harbor.

The mood was melancholy, or maybe that was my own. Once you've seen a place transformed by euphoria
and shared experience (and a lot of shared bar tabs), it can bring you down hard to watch it revert
to calm and the customary facades, to a more individualistic way of getting through the day.

The party here is not quite over. The Paralympic Winter Games start March 12. But the big bash has
ended, the one that required nearly a decade to plan and about two weeks to experience; the one that
cost billions and allowed Vancouver to raise its profile -- for worse, then better -- on a daily
basis.

Though bills are now due, the collective memories are still there, in heads and hearts. Before they
fade, this is the time to think about where the Winter Olympics go from here.

A few reminders to future host cities and the International Olympic Committee members who choose
them:

Weather matters: Winter sports -- at least the outdoor species -- are already vulnerable enough to
postponements without having to deal with rain or springtime temperatures. Global weirding means the
number of places with truly dependable Winter Olympic weather will probably shrink, not increase.

In defense of Cypress Mountain, where snow had to be trucked in, Vancouver did have its warmest
January in more than 100 years of record keeping. But the drizzle and slush were not restricted to
the northern suburbs. Whistler, the real mountain venue, also had plenty of mild-temperature issues,
particularly at lower altitudes, which was not a surprise to those who have attended World Cup
events there in the past.

Watching a halfpipe competition or slalom run take place in the rain is a buzz killer, just as it
was in Nagano, Japan, in 1998. Above all, the athletes deserve better at the biggest event of their
careers.

It was encouraging -- even if it came a bit late -- to hear Jacques Rogge, the International Olympic
Committee president, confirm on Sunday that evaluation committees will be asked to take a closer
look at weather data when they examine future candidates. But it needs to be a determining factor.
Meanwhile, here's hoping the Olympians get more snow than rain in 2014 in and around Sochi, Russia,
a city on the Black Sea much better known as a summer resort.

Winter sports culture matters: Crowds were generally enthusiastic at these Olympics, even if
Canadians had plenty to learn about biathlon and ski jumping. But the memory of the flat atmosphere
and sparse crowds in the mountain venues near (but not that near) Turin, Italy, in 2006 remains
strong.

The Winter Olympics lack the universality of the summer version. With the current dates in February,
no Southern Hemisphere nation could realistically stage them. There are far fewer countries whose
athletes make an impact and far fewer countries where winter sports are genuinely popular.

At the 2008 Summer Olympics in Beijing, 55 nations won gold medals and 87 won medals. In Vancouver,
19 nations won gold medals and just 25 won medals, one fewer than in Turin.

The temptation for the I.O.C. is to try to use the Winter Games as a developmental lever, which was
part of the thinking behind awarding them to Sochi. Russia, despite its famously harsh winters,
lacks quality Alpine ski resorts and other elite winter sports facilities.

But this promotional philosophy should be used sparingly. The Winter Games are best when the public
understands or connects deeply with the events, as it did in Lillehammer, Norway, in 1994 or, with
different priorities, in Salt Lake City, Utah, in 2002. The principle of continental rotation that
applies to the Summer Games should not be the rule in the winter. The I.O.C. should put more
emphasis on the best venue, even if it means repeats. And it should let its winter federations, like
the International Ski Federation and International Biathlon Union, do the global missionary work by
moving around their world championships instead.

The Olympics are too important -- both to the athletes and their public -- to be a testing ground.

Green matters: The Olympics, despite the challenges posed by their growth, have made progress in
being environmentally friendly. But there are big frustrations surfacing with the construction-heavy
Sochi Games. One of the biggest issues in modern Winter Games remains the sliding tracks. The one in
Cesana Pariol, built on pristine land for the Turin Games, is already in danger of being closed down
because of funding problems (the nearby ski jumps are also little used). There must be a better way
than to build expensive, environmentally invasive facilities in places with no culture of sliding
sports.

Building an Olympic track makes sense in countries without one like Japan or Norway. But Italy
already had a track built for the Olympics in Cortina d'Ampezzo. Canada already had one in Calgary,
Alberta, built for the 1988 Games. The new track in Whistler now has a questionable future after the
death of Nodar Kumaritashvili in training and concerns about its speed.

Better to use existing sites than risk another white elephant, or maybe it's better to hire
specialists from Tyrol who rebuild the naturally refrigerated track in St. Moritz, Switzerland, each
winter.

Unity mattersThe Winter Olympics have become all about increasingly distant clusters, with ice
sports in the city and snow sports and sliding sports in the mountains. Even with much reduced
traffic, it was about two hours and sometimes more between Vancouver and the Whistler venues --
dissuasive for spectators and for athletes.

The dual concept makes sense if you buy into the theory that you need increasingly substantial host
cities to support the Olympic supporting cast. But something precious is lost, and this is far from
the Olympic ideal of assembling the youth of the world every four years if all they do is assemble
for the opening and closing ceremonies and otherwise remain in their separate Olympic villages.

Sochi, in theory, should be less far-flung, with the resort of Krasnaya Polyana, site of the snow
sports, lying about 80 kilometers, or 50 miles, from the city. Two of the candidates for the 2018
Games -- Munich and Pyeongchang, South Korea -- are also emphasizing compact plans. Pyeongchang's
organizers claim that all venues would be within 30 minutes of each other.

That's the kind of clustering the Winter Olympics needs.

\section{Ads Posted on Facebook Strike Some as Off-Key}

\lettrine{F}{acebook}\mycalendar{Mar.'10}{05}, the world's biggest social network, is selling more
ad spots to big companies like Wal-Mart Stores, Procter \& Gamble and PepsiCo.

But the site's pages are also home to countless ads from smaller companies that can be funny, weird
or just plain creepy -- those suggesting you are, say, eligible to get a free iPad because you are
exactly 26 years old, or entreaties to see what your offspring would look like if you had a child
with a celebrity.

Odd Web ads, like the dancing women promoting mortgage brokers, are not new. But on social networks
like Facebook, where people go to communicate with one another, advertisers seem to be trying
especially hard to intrude on the conversation.

The so-called self-service ads on the site, from the likes of video game start-ups, herbal
supplement makers, sweepstakes companies and wedding photographers, are shown on the right side of
most pages. Many advertisers who use the self-service system are tempted to go as far as possible in
making ads that attract attention and appear relevant, aided by the information that people give to
Facebook.

``When it works, it's amazingly impactful, but when it doesn't work, it's not only creepy but
off-putting,'' said Tim Hanlon, a principal at the consulting firm Riverview Lane Associates of
Chicago. ``What a marketer might think is endearing, by knowing a little bit about you, actually
crosses the line pretty easily.''

One campaign that flooded the site in recent weeks, before Facebook cracked down on it, tries to
take advantage of consumer interest in Apple's iPad. ``Are you a fan of Eddie Izzard? We need 100
music and movie lovers to test and KEEP the new Apple iPad,'' one version of the ad says. Louis
Allred Jr., 29, a Facebook user in Los Angeles who saw the ad, said he figured it was shown to him
because he or a friend had expressed enthusiasm for Mr.~Izzard, a British comedian, on their
profiles.

``It doesn't seem like they are using the information in sensible ways,'' Mr.~Allred said.

Mr.~Allred was also skeptical about the test-an-iPad offer. The ad sends people to
Prize-Rewards.net, a site that appears to be based in Vancouver, Canada, and tries to get people to
pursue rewards by signing up for memberships in services like Netflix. Efforts to reach the company
were unsuccessful.

Dan Rose, vice president for business development at Facebook, compared the company's self-service
ad platform to the early versions of Google's highly profitable AdWords system. He predicted that
the quality of the promotional messages on the site would improve as more companies began to use it.

``A year ago, we had lower-quality ads, and a year from now we will have higher-quality ads,''
Mr.~Rose said. ``It's early, but we have made a lot of progress.''

The self-service ad system offers any person or company the opportunity to quickly design an ad and
aim it toward finely segmented groups of Facebook's 400 million members, based on gender, age,
location and preferences like favorite movies and activities.

For example, a promoter can advertise tickets to a band's concert to the select group of Facebook
users who live in the area and have mentioned that band on their profile page or status updates. Or
a wedding photographer can show ads only to people in a certain city who have switched their
relationship status to ``engaged.''

``I've had good success with it on a limited budget,'' said Carter Rose, a wedding photographer in
Dallas who has paid \$1,700 for Facebook ads over the last few months and as a result booked three
weddings at \$3,500 each.

Many advertisers view a site like Facebook as a good place for playful messages. Cordarounds, a San
Francisco pants company, talks about ``pigeon matadors'' and ``trouserbots'' that make pants in a
``zero-G environment'' to capture attention for its products on Facebook.

``It's a conversational medium. If you can be the originator of that conversation, you are fitting
in,'' said Chris Lindland, the company's chief executive, who uses Facebook to advertise solely to
young men in the 20 largest urban areas.

From the perspective of many users, the tailored ads can often seem, at best, presumptuous.

Women who change their status to ``engaged'' on Facebook to share the news with their friends, for
example, report seeing a flood of advertisements for services and products like wedding
photographers, skin treatments and weight-loss regimens.

Then there are the nonsensical ads, which highlight the fact that Facebook does not have employees
to review each ad, but relies largely on member feedback to flag inappropriate messages. That allows
ads to slip through that run afoul of the company's policies. ``Laws now allow U.S.~residents to
reduce their credit card debt by more than 50 percent,'' reads one ad for a debt-relief Web site,
which is accompanied by a surreal image of four tiny babies lying in the palm of a hand.

``It caught my eye, but it sure didn't win my trust,'' said Paige Kobert, 32, an advertising
employee from San Francisco who saw the ad recently. ``It just seemed like a scam.''

A Facebook spokesman, Brandon McCormick, said ads like that one and the Eddie Izzard iPad ad
violated a Facebook policy that requires the text and photo in the ad to be related to what is being
advertised.

Mr.~McCormick said that in the last few weeks, Facebook had begun enforcing that provision after
users complained about the flood of ads aimed at alumni from a certain university or people of a
certain age (``Hey, 32-Year-Olds''), with messages that had nothing to do with those
characteristics.

Such ads are the equivalent of a direct mail ``special offer for John Smith'' that turns out to be
not so special.

Facebook users seem most perplexed when they are left wondering why they have received a customized
ad.

Jess Walker, 22, from central Florida, was recently presented an ad for Plan B, the morning-after
pill.

``What do I have on my Facebook page that would lead them to believe I would need that?'' she asked,
adding that she did not want her sexual behavior called into question.

Ms.~Walker noted that Facebook allowed her to delete the ad from the page she was viewing. -- which
would then make Facebook less likely to allow that particular ad to be sent to other people.

Mr.~Rose from Facebook said the ads on the site were becoming more professional and straightforward
as the company updated its policies and enforced them.

``When you have got the platform that we have, you are naturally going to attract people who are
going to game it,'' Mr.~Rose said.

\section{Obama Calls for `Up or Down Vote' on Health Care Bill}

\lettrine{P}{resident}\mycalendar{Mar.'10}{05} Obama, beginning his final push for a health care
overhaul, called Wednesday for Congress to allow an ``up or down vote'' on the measure, and sketched
out an ambitious -- and, some Democrats said, unrealistic -- timetable for his party to pass a bill
on its own within weeks.

\myphoto{Obama.Calls.for.Up.or.Down.Vote.on.Health.Care.Bill.jpg}{President Obama called on
Wednesday for an ``up or down vote'' on health care overhaul.}

``I believe the United States Congress owes the American people a final vote on health care
reform,'' Mr.~Obama said during a 20-minute speech in the East Room of the White House. He said
there was no point in starting over, as Republicans are demanding, and called on nervous Democrats
to stick with him, declaring there was no reason ``for those of us who were sent here to lead to
just walk away.''

The speech, less than a week after Mr.~Obama held a high-profile televised health care forum, will
usher in what White House officials say will be their last campaign to bring Washington's long and
contentious health care debate to a close -- with a bill-signing ceremony at the end.

On Thursday, Health and Human Services Secretary Kathleen Sebelius will meet at the White House with
insurance industry executives to spotlight unpopular rate increases; next week, Mr.~Obama will
travel to Missouri and Pennsylvania to stump for the health care bill.

In his remarks, the president refrained from using the word ``reconciliation,'' the parliamentary
tactic that Democrats are expected to employ to avoid a Republican filibuster and win passage with a
simple majority. But he made clear that was his intent, and reminded Americans that despite current
Republican objections, other major bills had been passed using the same tactic.

``Reform has already passed the House with a majority. It has already passed the Senate with a
supermajority of 60 votes,'' Mr.~Obama said. ``And now it deserves the same kind of up or down vote
that was cast on welfare reform, that was cast on the Children's Health Insurance Program, that was
used for Cobra health coverage for the unemployed and, by the way, for both Bush tax cuts -- all of
which had to pass Congress with nothing more than a simple majority.''

Republicans were furious.

``They're making a vigorous effort to try to jam this down the throats of the American people, who
don't want it,'' the Senate Republican leader, Mitch McConnell of Kentucky, told reporters after
Mr.~Obama's remarks. ``We think that's a policy mistake, and we think resorting to these kind of
tactics, to thumb your noses at the American people, is something that ought to be resisted.''

On Capitol Hill, the strategy could prove a heavy lift for the House speaker, Nancy Pelosi, and the
Senate majority leader, Harry Reid, who are now under intense pressure from the White House to
translate Mr.~Obama's wishes for a final bill into legislative language. Both leaders issued
statements Wednesday praising Mr.~Obama and vowing to press ahead. But, noticeably, neither publicly
committed to Mr.~Obama's timetable.

Privately, Senate leadership aides said Mr.~Obama's deadline could be difficult to meet. The
tentative plan is for the House to adopt the bill passed by the Senate, and for both chambers to use
reconciliation to pass a package of changes that would bridge gaps between the initial House and
Senate versions.

But the final language must still be sent to the nonpartisan Congressional Budget Office for
evaluation, a process that takes time. Many aspects of the legislation remain unresolved, and
rank-and-file Democrats in the House remain deeply uneasy over both the substance of the bill and
the process by which it would be adopted.

Ms.~Pelosi does not yet have the votes she needs to pass the legislation. She faces complex
negotiations with both the moderate and liberal wings of her party to come up with a package that
can pass the House without deviating so much from the existing Senate version that Mr.~Reid would
have trouble assembling a majority for the final vote in the Senate.

``I am not inclined to support the Senate version,'' said Representative Shelley Berkley, Democrat
of Nevada, who voted for the House bill in November. ``I would like something more concrete than a
promise. The Senate cannot promise its way out of a brown paper bag.''

As Democrats prepared for a final showdown with Republicans, other potential stumbling blocks
emerged. House Democrats from New York met Wednesday with Ms.~Pelosi to discuss their concern that
the emerging bill would shortchange their state on Medicaid and other issues.

``I am very, very disappointed and unhappy,'' said Representative Eliot L.~Engel, Democrat of New
York. ``The White House is taking us for granted, and they shouldn't.''

Supporters of abortion rights, in and out of Congress, said Wednesday that they were alarmed at the
prospect that lawmakers might impose new restrictions on insurance coverage of abortion in the push
to enact sweeping health legislation. Representative Jan Schakowsky, Democrat of Illinois, said
language restricting insurers' ability to cover abortions ``remains in the president's proposal, and
we are very concerned about that.''

Friday will mark one year since Mr.~Obama laid out his plans for a health care overhaul with a
high-profile forum at the White House, where he engaged in a lively debate with lawmakers of both
parties and executives from the insurance, hospital and pharmaceutical industries. On Wednesday, the
scene at the White House was far different.

Mr.~Obama spoke, without taking questions, to a group of sympathetic medical professionals, many of
them clad in white lab coats to provide a TV-friendly image. After 12 months of legislative
hearings, town hall meetings, speeches, polls and debates, Mr.~Obama was in the position of selling
not only the public, but his own party, on his top domestic priority.

``The American people want to know if it's still possible for Washington to look out for their
interests and their future,'' Mr.~Obama said. ``They are waiting for us to act. They are waiting for
us to lead. And as long as I hold this office, I intend to provide that leadership. I don't know how
this plays politically, but I know it's right.''

Seeking to reassure wavering Democrats that he would back them up, he pledged to do ``everything in
my power to make the case for reform.''

Moments after he finished speaking, the White House announced plans for him to visit Pennsylvania
and Missouri -- states that are home to vulnerable Democrats like Representative Jason Altmire of
Pennsylvania and Representative Ike Skelton of Missouri, who were among 39 Democrats to vote against
the health measure when it passed the House last year. If Mr.~Obama is to sign his legislation into
law, he is going to have to convert some of those no votes into yeses; traveling to a lawmaker's
home state could be one way to do that.

Senior advisers to Mr.~Obama are betting that the politics of health care will eventually turn in
the party's favor, if the president can actually sign a bill into law. The legislation includes
popular restrictions on the insurance industry; some, like a provision barring insurers from
discriminating against children on the basis of pre-existing conditions, would take effect quickly
-- a point noted by Mr.~Obama's press secretary, Robert Gibbs.

``The president has always subscribed to the notion that the politics will catch up,'' Mr.~Gibbs said.
\clearpage
\renewcommand\listfigurename{\textit{Table of Figures}}
{\footnotesize\textit{\listoffigures}}

\end{document}
